\chapter{Research Methodology}
\label{ch:methodology}
\newpage

\section{Research Design Overview}
\label{sec:design-overview}

This study is a quantitative empirical benchmark of how well code written by AI coding agents
follows object-oriented design principles, in Java and in Python. The unit of analysis is the pull request: one change an
agent submitted to a real project, whether or not it was merged. The object of measurement is
adherence to the five SOLID principles, assessed one principle at a time with deterministic tools
and then checked against an independent reading of the code.

The study has two parts. Chapter~\ref{ch:literature} reviews work on AI code generation, coding
agents and software quality measurement, and identifies the gap the thesis addresses. This chapter
describes the empirical design, and Chapter~\ref{ch:results} applies it to two corpora of agent pull
requests, one in Java and one in Python, drawn from the AIDev dataset \citep{li2025aidev}.

Four design decisions define the study. They are justified below and in Section~\ref{sec:validity}.

\begin{enumerate}
  \item \textbf{Real pull requests, not newly generated repositories.} Running agents to generate
  many new codebases would be costly and would produce artificial tasks. Measuring code the agents
  actually submitted to projects gives up control over the tasks but gains realism, and each pull
  request is a finished change that can be attributed to one agent.
  \item \textbf{All three outcomes.} A merged pull request has passed review, so a corpus of merged
  requests measures the combined output of the agent and its reviewers. The corpus therefore
  contains merged, rejected (closed without merging) and open requests, with rejected and open
  requests over-sampled relative to the dataset. This makes it possible to examine what review
  filters out.

  A rejected request is not necessarily a request refused for its quality. Only requests with an
  explicit request for changes (\texttt{changes\_requested}) or a completed review followed by
  closure (\texttt{reviewed\_then\_closed}) record a judgement about the code; most closed requests
  were superseded, went stale or were abandoned. The outcome comparison is therefore always
  reported by these strata.
  \item \textbf{Two languages, Java and Python.} SOLID describes object-oriented design, and one can
  argue that measuring it in a language that does not require objects says more about the language
  than about the agent. Java removes that objection: classes are compulsory, \texttt{interface} is a
  keyword, \texttt{@Override} is explicit and subtyping is nominal. The reverse objection is that a
  result found in Java may only describe Java. The study is therefore carried out in both languages.
  Because the framework is parameterised by language (Section~\ref{sec:language-availability}), the
  whole study runs on Java and on Python, where classes are optional and no compiler checks a
  contract. A result is called general only if both languages show it.
  \item \textbf{Deterministic detectors, validated by a reading judge.} Every reported rate is
  computed by a program that gives the same output on every run. A large language model is used only
  to read source code and label a validation sample; it never computes a reported rate. Every
  agreement figure is therefore agreement between a detector and the judge
  (Sections~\ref{sec:reading-judge} and~\ref{sec:construct-validity}).
\end{enumerate}

The design is cross-sectional for each pull request. For Open/Closed, Liskov Substitution,
Interface Segregation and Dependency Inversion it compares the repository immediately before and
immediately after the change (Section~\ref{sec:delta-harness}). For Single Responsibility it
evaluates the touched classes in the tree after the change. It does not follow an agent across a
sequence of changes.

\subsection{Research Paradigm}

The study takes a post-positivist position. It looks for measurable and reproducible signals of
design quality while accepting that no static tool fully captures a semantic property such as
responsibility or substitutability. Each measurement is treated as a proxy for the underlying
principle, and the gap between proxy and principle is quantified with precision and recall against
the reading judge.

\subsection{Study Type}

The study is observational. The agents worked on their own tasks in repositories this study does not
control, before anything was measured. No treatment was assigned, so the design supports comparison
and description but not causal claims (Table~\ref{tab:study-type}).

\begin{table}[htb]
\centering
\caption{Study type.}
\label{tab:study-type}
\begin{tabular}{@{}L{3.4cm}L{10.2cm}@{}}
\toprule
\textbf{Dimension} & \textbf{Description} \\
\midrule
Purpose & Comparison of agents and validation of the measurement instruments \\
Unit of analysis & Pull request, with class-level and module-level detail \\
Subjects & Three AI coding agents: Claude Code, GitHub Copilot, OpenAI Codex \\
Outcomes & Merged, rejected (closed without merging) and open, with rejected and open
over-sampled \\
Languages & Java and Python, with the same framework in both \\
Data source & AIDev dataset \citep{li2025aidev}, curated subset \\
Data collection & Automated static analysis of the source before and after each pull request \\
Analysis & Per-agent and per-outcome rates (pooled and averaged over repositories), a joint
regression model, and validation precision and recall \\
\bottomrule
\end{tabular}
\end{table}

\section{Research Questions}
\label{sec:rqs}

The study answers the two research questions of Section~\ref{sec:research-questions}.
Table~\ref{tab:rq-summary} summarises how each is answered, and the paragraphs below describe the
analyses behind it.

\paragraph{RQ1, adherence and differentiation.}
\emph{Do coding agents differ in how well their pull requests adhere to the five SOLID principles,
and how many of the design problems in the code they touch do they introduce themselves, in Java
and in Python?}

RQ1 is answered by four analyses. The first three are carried out on the Java corpus and on the
Python corpus, and the fourth compares the two.

\textbf{Differences between agents.} On a common corpus, with repository and change size controlled,
do agents leave different structural footprints, and on which principles?

\textbf{Merged and unmerged requests.} A corpus of merged pull requests shows what maintainers
accepted, which is not the same as what agents produced. If the code in rejected requests were
structurally worse, the merge decision would be filtering on design quality, and a study of merged
code would attribute that filtering to the agents. Including rejected and open requests makes the
comparison possible.

\textbf{Locus of risk.} A class or a method can be written in full by one pull request. A dependency
cycle, an interface and its implementers, or a base class and its override involve at least two units,
and one pull request rarely writes both sides. Single Responsibility and Open/Closed violations lie
inside one unit; Liskov Substitution, Interface Segregation and Dependency Inversion violations lie
in a relation between units (Figure~\ref{fig:intra-inter}). The analysis compares the violations
already present in the code a pull request touched with the violations the pull request introduced.
The answer indicates which kinds of design damage a single contribution can cause and which kinds
build up over many contributions.

\textbf{Both languages.} Each answer could depend on the language. Change size could dominate
because Java code is verbose, an agent difference could reflect how one agent writes Java, and
relation-level violations could be rare because the Java compiler rejects some of them. The second
language tests these possibilities. A second corpus is drawn from the same dataset for the same
agents and outcomes. The fact extractor and the thresholds are replaced; the detector logic,
validation procedure, judging rubric and statistical model are reused. An answer to RQ1 is called
general only if both languages show it.

\begin{figure}[htb]
\centering
% Figure: where a violation lives, and what one pull request can reach.
% Two panels of equal size. The gold dashed overlay is the pull request's diff.
% Left: the whole class sits inside the diff, so the change can create the
% violation. Right: the violation lies on the inheritance edge between two files;
% the diff covers only the subclass file, so the base side predates the change.
% Vertical plan per panel: title -0.32, subtitle -0.66, content -0.95 to -3.8,
% footnote from -3.95. Rendered at 150 dpi and checked for overlaps.
\begin{tikzpicture}[
  font=\scriptsize,
  panel/.style={draw=issatblue!30,fill=lightblue!25,rounded corners=6pt,line width=0.6pt},
  ttl/.style={font=\footnotesize\bfseries,text=issatblue,anchor=west},
  sub/.style={font=\tiny,text=accentgray,anchor=west},
  file/.style={draw=accentgray!80,fill=white,rounded corners=2pt,line width=0.6pt},
  fname/.style={font=\tiny\ttfamily,text=accentgray,anchor=north east,inner sep=2pt},
  code/.style={font=\tiny\ttfamily,anchor=west,inner sep=1pt},
  codeline/.style={fill=accentgray!30,rounded corners=0.5pt},
  diff/.style={draw=issatgold!85!black,dashed,line width=1.1pt,rounded corners=4pt,fill=issatgold!12},
  tab/.style={fill=issatgold!85!black,text=white,font=\tiny\bfseries,rounded corners=2pt,inner sep=2pt},
  foot/.style={font=\scriptsize\itshape,text=black!75,align=center,text width=7.0cm},
]
% ================================================================ left panel
\draw[panel] (0,0) rectangle (7.4,-4.85);
\node[ttl] at (0.25,-0.32) {Inside one unit};
\node[sub] at (0.25,-0.66) {Single Responsibility, Open/Closed};

\draw[diff] (0.7,-0.95) rectangle (6.7,-3.8);
\node[tab,anchor=north east] at (6.62,-1.02) {pull request diff};
\draw[file] (0.9,-1.35) rectangle (6.5,-3.6);
\node[fname] at (6.48,-1.37) {OrderService.java};
\node[code] at (1.05,-1.62) {class OrderService \{};
\foreach \y/\c/\m/\a/\b in {-2.0/srpcolor/parseOrder/2.3/1.6,
                             -2.5/issatgold/sendEmail/1.9/2.4,
                             -3.0/dipcolor/chargeCard/2.5/1.4}{
  \fill[\c] (1.15,\y-0.26) rectangle (1.24,\y+0.1);
  \node[code,text=\c!75!black] at (1.3,\y) {\m()};
  \fill[codeline] (1.36,\y-0.17) rectangle ({1.36+\a},\y-0.12);
  \fill[codeline] (1.36,\y-0.25) rectangle ({1.36+\b},\y-0.2);
}
\node[code] at (1.05,-3.4) {\}};

\node[foot,anchor=north] at (3.7,-3.93) {The change can write the whole class,\\so one pull request can create the violation.};

% =============================================================== right panel
\draw[panel] (7.9,0) rectangle (15.3,-4.85);
\node[ttl] at (8.15,-0.32) {Between units};
\node[sub] at (8.15,-0.66) {Liskov, Interface Segregation, Dependency Inversion};

% base file: not part of this change
\draw[file] (8.3,-1.0) rectangle (11.5,-2.05);
\node[fname] at (11.48,-1.02) {Shape.java};
\node[code] at (8.42,-1.4) {class Shape \{};
\node[code] at (8.55,-1.7) {area() \{ return w*h; \}};
\node[font=\tiny\itshape,text=accentgray,anchor=north] at (9.9,-2.12) {not in this pull request};

% subclass file: inside the diff
\draw[diff] (11.45,-2.35) rectangle (15.1,-3.84);
\node[tab,anchor=north east] at (15.02,-2.42) {pull request diff};
\draw[file] (11.65,-2.72) rectangle (14.9,-3.66);
\node[fname] at (14.88,-2.74) {Circle.java};
\node[code] at (11.77,-2.98) {class Circle};
\node[code] at (11.9,-3.2) {extends Shape \{};
\node[code,text=ispcolor!85!black] at (11.9,-3.44) {area() \{ throw ... \}};

% the relation, where the violation lives
\draw[-{Triangle[open,length=2.4mm,width=2.4mm]},accentgray!90,line width=0.9pt]
  (12.3,-2.72) -- (12.3,-1.52) -- (11.5,-1.52);
\node[circle,fill=ispcolor,text=white,font=\tiny\bfseries,inner sep=0pt,minimum size=0.34cm]
  at (12.3,-2.02) {!};
\node[font=\tiny,text=ispcolor!85!black,anchor=west,align=left] at (12.55,-1.72)
  {the violation lies\\on this relation};

\node[foot,anchor=north] at (11.6,-3.93) {Each side is usually written by a different change,\\so these violations build up over many changes.};
\end{tikzpicture}

\caption{Intra-unit and inter-unit violations.}
\label{fig:intra-inter}
\end{figure}

Figure~\ref{fig:intra-inter} shows the distinction behind the locus of risk. The dashed gold area is
what one pull request changes. On the left the whole class is inside the change, so the change can
create the violation. On the right the violation lies on the relation between two files, and the
change covers only one of them.

\paragraph{RQ2, measurability and construct validity.}
\emph{Can SOLID adherence be measured at repository scale with deterministic structural detectors,
and how well do these detectors agree with an independent reading judge and with established
analysers run on the same code?}

SOLID principles concern meaning and intent, while static metrics describe structure. RQ2 asks how
far a structural detector can be trusted, and answers with precision and recall per principle
against labels from a judge that reads the source, in both languages. For Java, a second check
compares the class-level detectors with three published analysers, CK, PMD and DesigniteJava, run on
the same source trees (Section~\ref{sec:validation-independent}). The two checks can fail in
different ways: the judge can err in the same direction as a detector, while an independently built
tool has no access to the detector's rules.

\begin{table}[htb]
\centering
\caption{Research questions and how each is answered.}
\label{tab:rq-summary}
\begin{tabular}{@{}cL{3.0cm}L{3.2cm}L{5.4cm}@{}}
\toprule
\textbf{RQ} & \textbf{Focus} & \textbf{Construct} & \textbf{Method} \\
\midrule
RQ1 & Agents, outcomes, locus of risk, two languages & SOLID adherence per principle & Rates per agent
and outcome, pooled and per repository; joint regression model; violations present compared with
violations introduced; 4{,}245 Java and 4{,}946 Python pull requests \\
RQ2 & Construct validity & Detector fidelity & Precision and weighted recall on 1{,}495 judged Java
and 980 judged Python entities, including a Java holdout sample; agreement of the Java detectors with
three external analysers \\
\bottomrule
\end{tabular}
\end{table}

\section{Research Strategy}
\label{sec:strategy}

\subsection{A Benchmark over Real Pull Requests}

A fixed corpus of agent pull requests is processed by one deterministic analysis procedure, and the
results are compared across agents. Mining finished work has two advantages: the code was written
for a real task in a real project, and each pull request can be attributed to one agent. The cost is
that task difficulty cannot be held constant across agents. Repository controls
(Section~\ref{sec:repo-mix}) and the joint model reduce this problem, and
Section~\ref{sec:limitations} states what remains.

The corpus comes from the AIDev-7.6M release, whose master tables hold 7,685,281 agent pull
requests in 957,209 repositories. The agents are very unevenly represented: OpenAI Codex has
3,844,602 pull requests and Devin 77,791. The working corpus is therefore sampled by agent and
outcome; Section~\ref{sec:dataset} describes the curation.

\subsection{One Deterministic Detector per Principle}

Each principle is implemented as a pure function over extracted facts. The function receives the
facts for a source tree and returns a structural measurement and a flag. It does not depend on how
the source was obtained or on any external service, so the same input always gives the same output.
The Single Responsibility detector runs on the touched classes, and the other four run on complete
before and after trees through the same interface. Section~\ref{sec:proxy-mapping} summarises each
detector and Chapter~\ref{ch:results} gives the exact rules and thresholds.

\subsection{The Before/After Delta}
\label{sec:delta-harness}

For Open/Closed, Liskov Substitution, Interface Segregation and Dependency Inversion, a single
snapshot of the repository cannot show what the agent changed. The analysis therefore reconstructs
the repository immediately before and immediately after the pull request, analyses both trees, and
charges the pull request only with the difference: a new dependency cycle, a newly widened interface
or a newly incompatible override. For a merged request the two trees
come from the merge commit and its first parent; for an unmerged request they come from the
request's branch point and head (Section~\ref{sec:dataset}).

Figure~\ref{fig:delta-harness} illustrates the delta for Dependency Inversion. The same analysis runs
on the tree before and after a pull request. In the example the pull request edits modules \texttt{a}
and \texttt{c} (gold rings) and adds the import \texttt{c}$\rightarrow$\texttt{a}, which makes three
module pairs mutually reachable; only this difference is charged to the pull request.

\begin{figure}[htb]
\centering
% Figure: the before/after delta for Dependency Inversion.
% Left: the import graph before the pull request. Right: after it. The pull
% request touches modules a and c (gold rings) and adds one import, c -> a,
% which closes the cycle a -> b -> c -> a. Bottom: the subtraction the detector
% performs, and what is charged. Module positions are the same in both panels
% so the reader compares like with like; labels sit outside every edge.
\begin{tikzpicture}[
  font=\scriptsize,
  panel/.style={draw=issatblue!30,fill=lightblue!25,rounded corners=6pt,line width=0.6pt},
  ttl/.style={font=\footnotesize\bfseries,text=issatblue,anchor=west},
  sub/.style={font=\tiny,text=accentgray,anchor=west},
  mod/.style={circle,draw=accentgray,fill=white,line width=0.7pt,minimum size=0.78cm,
              inner sep=0pt,font=\scriptsize\ttfamily},
  touched/.style={mod,draw=issatgold!90!black,line width=1.6pt},
  imp/.style={-{Latex[length=2.2mm]},draw=accentgray,line width=0.8pt},
  newimp/.style={-{Latex[length=2.6mm]},draw=ispcolor,line width=1.5pt},
  cyc/.style={fill=ispcolor!9,draw=ispcolor!45,dashed,line width=0.6pt,rounded corners=12pt},
  chip/.style={draw=#1,fill=#1!10,rounded corners=3pt,inner sep=4pt,align=center,font=\scriptsize},
]
% ---------------------------------------------------------------- panels
\draw[panel] (0,0) rectangle (6.4,-4.2);
\node[ttl] at (0.25,-0.32) {Before};
\node[sub] at (0.25,-0.64) {tree the pull request started from};

\draw[panel] (8.9,0) rectangle (15.3,-4.2);
\node[ttl] at (9.15,-0.32) {After};
\node[sub] at (9.15,-0.64) {tree the pull request produced};

% ------------------------------------------------------------ before graph
\begin{scope}[shift={(0,0)}]
  \node[mod] (a1) at (1.3,-1.55) {a};
  \node[mod] (b1) at (4.9,-1.55) {b};
  \node[mod] (c1) at (3.1,-3.35) {c};
  \node[mod] (d1) at (5.4,-3.35) {d};
  \draw[imp] (a1) -- (b1);
  \draw[imp] (b1) -- (c1);
  \draw[imp] (b1) -- (d1);
\end{scope}

% ------------------------------------------------------------- after graph
\begin{scope}[shift={(8.9,0)}]
  \path[cyc] (0.72,-1.02) -- (5.5,-1.02) -- (3.1,-3.98) -- cycle;
  \node[font=\tiny\itshape,text=ispcolor!80!black,anchor=north] at (3.1,-1.1) {new cycle};
  \node[touched] (a2) at (1.3,-1.55) {a};
  \node[mod] (b2) at (4.9,-1.55) {b};
  \node[touched] (c2) at (3.1,-3.35) {c};
  \node[mod] (d2) at (5.4,-3.35) {d};
  \draw[imp] (a2) -- (b2);
  \draw[imp] (b2) -- (c2);
  \draw[imp] (b2) -- (d2);
  \draw[newimp] (c2) -- (a2);
  \node[font=\tiny\bfseries,text=ispcolor,anchor=east] at (1.95,-2.62) {new import};
\end{scope}

% --------------------------------------------------------- the pull request
\draw[-{Latex[length=3mm]},issatgold!90!black,line width=2pt] (6.65,-2.1) -- (8.65,-2.1);
\node[font=\scriptsize\bfseries,text=issatgold!80!black,anchor=south] at (7.65,-2.05) {pull request};
\node[font=\tiny,text=accentgray,anchor=north,align=center] at (7.65,-2.2) {edits \texttt{a}, \texttt{c}};

% -------------------------------------------------------------- the result
\node[chip=accentgray,anchor=north west,text width=4.3cm] (sub1) at (0,-4.5)
  {mutually reachable pairs before:\\ \texttt{none}};
\node[font=\small,text=accentgray] (minus) at (5.1,-4.95) {$\rightarrow$};
\node[chip=ispcolor,anchor=north west,text width=4.6cm] (res) at (5.55,-4.5)
  {new pairs after: \texttt{(a,b) (b,c) (a,c)}\\all involve a touched module};
\node[font=\small,text=accentgray] at (10.9,-4.95) {$\rightarrow$};
\node[chip=issatgold!90!black,anchor=north west,text width=3.6cm,font=\scriptsize\bfseries] at (11.35,-4.5)
  {flag charged to this\\pull request};
\end{tikzpicture}

\caption{The before/after delta for Dependency Inversion.}
\label{fig:delta-harness}
\end{figure}

Two attribution rules change what is measured. First, classes an agent created are scored separately
from existing classes it edited. Scoring only the touched classes would charge the agent with the
debt already in the repository (Section~\ref{sec:findings}, Finding 4). Second, for the
relation-level principles a violation is charged when either side of the relation is in a touched
file. An agent that edits a base class or widens an interface is therefore charged for the
subclasses and implementers it breaks in files it did not open.

For ISP, LSP and most DIP cases, a single pull request rarely creates a whole violation, because the
violation spans code the request did not write. Each of these detectors therefore also records an
inventory of the relevant structures at the head of the pull request (every interface, override
pair or dependency edge), reported separately from the per-request delta.

\subsection{Validation by a Reading Judge}
\label{sec:reading-judge}

For each principle a stratified sample of flagged and unflagged items is drawn. A large language
model of the Claude family acts as the judge: it reads the source of each item and labels it as a
violation or not, with a written reason. For LSP the judge sees the overriding method next to its
base method. The detector's flag is then scored against these labels, so every precision and recall
figure in this thesis is agreement between the detector and the judge. The label columns are named
accordingly (for example \texttt{llm\_srp\_risk}). The judge sees the entity and its source but not
the detector's verdict.

\subsection{Screens, Suppressors and Refinements}

The Single Responsibility and Open/Closed detectors are designed as screens: their thresholds aim to
miss few violations and accept false alarms. Each has a suppressor that removes one named class of
false positives. For Single Responsibility, a class whose methods share a common vocabulary
(conceptual cohesion, C3) is not flagged. For Open/Closed, a dispatch over an enumeration is not
flagged, because the language fixes its set of values. Whether the screens reach high recall is
measured in Section~\ref{sec:instrument-validation}.

A change to a detector was accepted only if the reason it should work could be stated in advance.
A restriction derived from reading false positives can be justified by the rows it removes. A
threshold found by trying values until the score improves cannot be distinguished from fitting
noise.

Two refinements reported in Chapter~\ref{ch:results} were derived by reading the false positives of
a first sample. Scoring them on the same sample would overstate their performance, so each was
scored again on a holdout: 200 further entities drawn with a different seed, excluding every row
judged before, and judged before the two rounds were compared.

\subsection{Controlling for Repository Mix}
\label{sec:repo-mix}

Because the corpus is mined, an agent's pooled rate can be dominated by one large repository in which
it happens to work often. Every comparative rate is therefore reported twice: pooled over all of an
agent's units, and as the mean of per-repository rates (the macro average). The share of an agent's
pull requests that come from its largest repository is reported as a concentration check. Where the
two rates disagree, the macro average is preferred.

\subsection{Language Selection}
\label{sec:language-availability}

Three detectors work on classes: the God Class rubric (SRP), the interface width and stub rule (ISP)
and the override comparison (LSP). They can only measure an agent if the corpus contains enough
classes.

Java meets this condition by design. Every Java compilation unit declares a type, so every pull
request that changes Java source touches a class. Each principle also has a direct counterpart in
the syntax: \texttt{interface} declares a contract, \texttt{@Override} marks an override, subtyping
is nominal, and the compiler rejects several violations that other languages allow at run time.

The last property has a consequence developed in Chapter~\ref{ch:results}. Three of the signals a
Liskov detector needs in a dynamically typed language cannot occur in Java code that compiles: a
changed parameter list makes an overload, parameter types are invariant, and visibility cannot be
reduced. The Java detector keeps only the signals the compiler does not check.

Only the fact extractor depends on the language. It turns source code into a language-neutral
record of types, members, supertypes, imports and member accesses, and every detector, calibration
step, validation procedure and report reads that record. A second language therefore needs a new
extractor, not a new study.

Python was chosen as the second language because it lacks the properties that made Java convenient.
A Python module does not need to declare a class; a contract can be an \texttt{abc.ABC}, a
\texttt{typing.Protocol} or a base class whose methods raise \texttt{NotImplementedError}; and
nothing checks an override against its base before the program runs. The class precondition is met
in practice: the analysable Python pull requests touch 81{,}149 production classes. The compiler
argument also makes a testable prediction. If \texttt{javac} removes a family of Liskov faults, those
faults should make up a large share of the Liskov violations a judge confirms in Python.
Section~\ref{sec:python-validation} reports that they do.

\section{Mapping SOLID to Structural Detectors}
\label{sec:proxy-mapping}

This section describes, at the level of method, how each principle is turned into a computable
detector. Chapter~\ref{ch:results} gives the exact thresholds, and Table~\ref{tab:proxy-map}
summarises the mapping.

\begin{description}[leftmargin=1.4cm,style=nextline]
\item[SRP (Single Responsibility).] The targeted smell is the God Class, a class doing several
unrelated jobs. The detector combines a six-metric rubric of size, complexity and cohesion (class
lines of code, method count, weighted methods, per-method complexity and parameters, and LCOM) with
conceptual cohesion, C3 \citep{marcus2008c3}. A class is flagged when it crosses at least two
thresholds and its methods do not share a common vocabulary. Size alone cannot separate a large
class with one topic from a large class with several, which is why cohesion is combined with size.

\item[OCP (Open/Closed).] The targeted smell is a type or tag switch that must be edited whenever a
new variant is added. The detector finds dispatch chains over the type or tag of one subject
(\texttt{instanceof} ladders, \texttt{getClass()} comparisons and \texttt{switch} statements on an
enumeration or a tag value), flags chains longer than a calibrated number of branches, and
suppresses chains over a set of values fixed by the language. Comparing the before and after trees
shows whether a pull request extended such a chain or added a new one.

\item[LSP (Liskov Substitution).] The targeted smell is an override that cannot stand in for its
base: an inherited method that is refused, a precondition strengthened by a new argument check, or a
postcondition weakened by returning \texttt{null} where the base never does. In Java the compiler
already rejects signature faults (a changed parameter list is an overload, parameter types are
invariant, and visibility cannot be reduced), so the Java detector looks only for faults the compiler
does not check. It works over the resolved type hierarchy, which avoids building each of the 1{,}010
repositories. The Python detector adds the signature signals
(Section~\ref{sec:python-instrument}).

\item[ISP (Interface Segregation).] The targeted smell is a wide interface whose implementers must
stub methods they do not need. The detector records the width of every contract (its number of
abstract members) and, for every implementer, the share of members it stubs. A pull request is
flagged when it creates or widens a contract beyond a calibrated width, or leaves an implementer
with forced stubs. Width is a screen: a wide interface whose implementers use every member is not a
violation. The forced stub is direct evidence.

\item[DIP (Dependency Inversion).] The targeted smell is a module that depends on concrete details,
and in particular a dependency cycle. The detector builds the internal import graph before and
after the pull request, finds strongly connected components with Tarjan's algorithm
\citep{tarjan1972scc}, and flags a change that makes two modules newly reachable from each other. It
also records the change in mean distance from Martin's main sequence and whether new dependencies
point to abstract or concrete modules, as description only. The distance signal confirmed none of
its flags when validated, so it does not enter the flag (Section~\ref{sec:zones-detail}).

\item[Component zones (descriptive only).] A sixth evaluator places every package on Martin's
abstractness and instability plane \citep{martin2003agile}: the zone of pain (concrete and
stable), the zone of uselessness (abstract and unused) and the main sequence between them. It uses
the $A$, $I$ and $D$ values the DIP detector computes and raises no flags. A package with no classes
or no internal dependencies has undefined coordinates and is reported as unplaceable, not as zero;
recording zero would place every class-free utility package in the zone of pain.
\end{description}

\begin{table}[htb]
\centering
\caption{From principle to detector (details in Chapter~\ref{ch:results}).}
\label{tab:proxy-map}
\small
\begin{tabular}{@{}L{1.75cm}L{2.6cm}L{4.1cm}L{4.35cm}@{}}
\toprule
\textbf{Principle} & \textbf{Smell} & \textbf{Detector} & \textbf{Main metric or rule} \\
\midrule
SRP & God Class & size and complexity rubric with a cohesion suppressor & LCOM; C3 conceptual cohesion \\
OCP & type or tag switch & dispatch-chain length with a closed-set suppressor & branch count;
closed-set test \\
LSP & incompatible override & faults the compiler does not check, over the resolved hierarchy &
refused bequest; argument checks; returning null \\
ISP & wide interface, forced stubs & contract width and implementer stub ratio & abstract-member
count; stub ratio \\
DIP & dependency cycle & import graph before and after & new cyclic module pairs (Tarjan) \\
Zones & descriptive only & package position on the $A$/$I$ plane & pain, main sequence,
uselessness, intermediate, unplaceable \\
\bottomrule
\end{tabular}
\end{table}

\section{Validity}
\label{sec:validity}

This section lists the design decisions that protect validity. The threats that remain are
collected in Section~\ref{sec:limitations}.

\paragraph{Internal validity.} Every measurement is produced by a deterministic program from source
text, so every reported number can be regenerated from the corpus. Thresholds are calibrated from
percentiles of each corpus before any labelling, so a detector is never tuned on the sample used to
score it. A violation is charged to a pull request only through a before and after comparison, so a
violation that already existed is not charged to an agent that edited nearby code.

\paragraph{Construct validity.} Each detector is validated on a stratified sample of flagged and
unflagged entities. A sample of flagged entities alone can estimate precision but cannot find missed
violations. Every label is written by a judge that reads the entity's source and records a reason.
For Java, the class-level detectors are also compared with three independently developed analysers
(Section~\ref{sec:validation-independent}).

\paragraph{Statistical validity.} Counts are reported next to rates, and proportions carry Wilson
score intervals, which behave well for small counts. Rates are reported pooled and as macro
averages over repositories. The joint model estimates agent effects with standard errors clustered
by repository, so many pull requests from one repository are not treated as independent. Because
several agent coefficients are tested, each is also compared with a Bonferroni threshold over the
agent coefficients of its language (Section~\ref{sec:agent-differentiation}).

\paragraph{External validity.} A result from one language may describe the language and not the
agents. The study addresses this by repeating corpus construction, calibration, validation and
modelling in Python with the same agents (Section~\ref{sec:python}). Each language's thresholds are
calibrated on its own corpus before its labels exist.

\paragraph{Reliability.} The pipeline can be resumed after interruption, facts are cached by the
content hash of the source tree, and each pure module has a self-check that runs as a script. A
mechanical audit checks that every judged row traces to the source file the judge was shown.

\section{Ethical Considerations}
\label{sec:ethics}

The study involves no human participants and no personal or proprietary data. Pull request metadata
comes from the public AIDev dataset. The repositories were public and carried a permissive licence
at the time of collection; they were cloned locally for analysis and are not redistributed. No model
was trained. The only use of a language model is as a reading judge on public source code during
validation.
